\section{Cloud Architektur}
\subsection{Cloud Native}
Cloud Native ist das erste mal 2013 angesprochen worden, von Netflix auf einen AWS reinvent talk. Es gibt keine Klare Defention für Cloud Native Architektur. Grundsatzlich soll Cloud Native aber dabei helfen Web Anwendung zu entwicklen die Skalierbar und Hochverfügbar sind. Gleichzeitg soll durch Cloud Native Agil entwickelt werden also eine Schnellst mögliche feature anpassung mit möglichst kleine Risiko. Dabei helfen Micoservice statt einer Monolithischen Architektur, Container, Dev Ops, CI/CD Pipilines. Auch Viele Lösungen wie Komponenten, Logging, Tracing sind Stadertsiert. 

\begin{figure}[H]
    \centering
    \includegraphics[width=\textwidth]{CloudNativeMap.png}
    \caption{Kubernetes Features}
\end{figure}

\subsubsection{DevOps}
Kombination aus \textit{Development} und \textit{Operations}.
Entwicklung und Betrieb arbeiten zusammen statt getrennt.
Kern-Ideen: Automatisierung von Builds, Tests und Deployments,
schnelle Feedbackschleifen (kurze Release-Zyklen) und gemeinsame
Verantwortung für Code und Infrastruktur.

\subsubsection{CI/CD -- Continuous Integration / Continuous Delivery}
\textbf{CI:} Jeder Code-Push wird automatisch gebaut und getestet,
Fehler werden sofort sichtbar.
\textbf{CD:} Nach erfolgreichem CI wird der Build automatisch
in eine Umgebung ausgeliefert (Staging oder Produktion).

\subsubsection{Microservices}
Architekturstil, bei dem eine Anwendung in viele kleine, unabhängige
Dienste aufgeteilt wird. Jeder Service hat eine klar abgegrenzte Aufgabe,
kommuniziert über APIs (REST oder Message Queues) und kann separat
deployed und skaliert werden. Gegenteil: \textit{Monolith}.

\begin{figure}[H]
    \centering
    \includegraphics[width=\textwidth]{CloudNativeBesipel.png}
    \caption{Kubernetes Features}
\end{figure}

\subsection{Architektur-Grundlagen}

\subsubsection{Warum wichtig?}
Pläne beeinflussen spätere Entscheidungen. Viele Änderungsanfragen
sind vermeidbar, sonst entstehen technische Schulden. Framework-Änderungen
allein lösen keine Grundprobleme, und eingeschränkte Komponentenwahl
wirkt sich auf Folgeprojekte aus (z.\,B. Lizenzen).

\subsubsection{Technische Schulden}
\textbf{Definition:} Konsequenzen bei schlechter Umsetzung.\\[4pt]
\textbf{Arten:}
\begin{itemize}
    \item Bewusste Schulden (strategische Entscheidungen)
    \item Unbeabsichtigte Schulden durch mangelnde Dokumentation,
          fehlende Automatisierung, fehlende Coding Standards
          oder Missachtung von Code-Hinweisen
\end{itemize}

\subsubsection{Ziele einer Lösungs-Architektur}
\begin{itemize}
    \item Skalierbarkeit, Resilienz, Sicherheit, Compliance
    \item Globaler Kundendienst mit gleicher Performance
    \item Wartbarkeit über Jahrzehnte
\end{itemize}

\subsubsection{Architektur-Perspektiven}
\begin{itemize}
    \item \textbf{Geschäftssicht:} Ziele, Vision, Stakeholder,
          Geschäftsanforderungen, Erfolgsmetriken
    \item \textbf{Prozesssicht:} Abläufe von Kunde zu Server,
          z.\,B. als Sequenzdiagramm oder BPMN
    \item \textbf{Konzeptionelle Sicht:} Architektur ohne
          spezifische Cloud-Provider-Komponenten
    \item \textbf{Lösungssicht:} Cloud-Service-Architektur,
          User Flow, ER-Diagramm
    \item \textbf{Sicherheitssicht:} Zugansverwaltung,
          Compliance-Anforderungen, Firewall
\end{itemize}

\newpage

\subsubsection{Architektur-Frameworks: Governance}
Dienen der unternehmensweiten Steuerung, Compliance-Nachweisen
und Risikokontrolle über vordefinierte Prinzipien, Gremien und Reviews.
Beispiele: TOGAF, ISO/IEC/IEEE 42010:2022, Zachman-Framework.\\[4pt]
\textbf{TOGAF} (The Open Group Architecture Framework): Framework zur
Entwicklung und Verwaltung von Unternehmensarchitekturen (People,
Processes, Technology) mit vier Domänen:
\begin{itemize}
    \item Geschäftsarchitektur
    \item Informationssystemarchitektur (Anwendungs- und Datenarchitektur)
    \item Technologiearchitektur
\end{itemize}

\subsubsection{Architektur-Frameworks: Praxisorientiert}
Stellen schnelle, verständliche Diagramme für Entwickler, DevOps und
Produktteams bereit. Zielen eher auf einzelne Lösungen und Applikationen
ab. Beispiele: C4-Modell, 4+1-Modell, PlantUML.\\[4pt]
\textbf{C4-Modell:} Unterstützt Software-Entwicklungsteams, ihre
Architektur bei Planung und Durchführung zu erklären.
4 Ebenen: Software System, Container, Component, Code.\\[4pt]
\textbf{Software Architecture Diagramming Maturity Model:}
\begin{itemize}
    \item \textbf{Level 1 -- Initial:} Keine Architekturdiagramme
    \item \textbf{Level 2 -- Ad hoc:} Diagramme mit ad-hoc-Abstraktionen
          in einem allgemeinen Diagramm-Tool
    \item \textbf{Level 3 -- Defined:} Diagramme mit definierten
          Abstraktionen und Notation, allgemeines Tool
    \item \textbf{Level 4 -- Modelled:} Definierte Abstraktionen in
          einem Modellierungs-Tool, manuell erstellt
    \item \textbf{Level 5 -- Optimising:} Modellelemente werden
          teamübergreifend geteilt, als querybare Datensätze genutzt
          und automatisch aus Sourcecode bzw. Deployment-Umgebung
          generiert
\end{itemize}

\subsection{TOGAF -- Architecture Development Method (ADM)}
\begin{figure}[H]
    \centering
    \includegraphics[width=\textwidth]{TOGAF.png}
    \caption{Kubernetes Features}
\end{figure}

Der ADM-Zyklus beschreibt, wie eine Unternehmensarchitektur schrittweise
entwickelt wird. Im Zentrum stehen immer die \textbf{Requirements} (Anforderungen),
auf die jede Phase kontinuierlich einzahlt.

\begin{itemize}
    \item \textbf{Preliminary:} Framework und Prinzipien festlegen
          (Vorbereitung, noch kein offizieller Buchstabe)
    \item \textbf{A -- Architecture Vision:} Ziele, Scope und
          Stakeholder-Erwartungen definieren
    \item \textbf{B -- Business Architecture:} Geschäftsprozesse
          und -strukturen modellieren
    \item \textbf{C -- Information System Architecture:}
          Anwendungs- und Datenarchitektur beschreiben
    \item \textbf{D -- Technology Architecture:} Infrastruktur,
          Plattformen und Technologien festlegen
    \item \textbf{E -- Opportunities \& Solutions:} Lücken zwischen
          Ist- und Soll-Zustand identifizieren, Lösungsoptionen bewerten
    \item \textbf{F -- Migration Planning:} Roadmap und Reihenfolge
          der Umsetzung planen
    \item \textbf{G -- Implementation Governance:} Umsetzung
          begleiten und sicherstellen, dass Architekturvorgaben
          eingehalten werden
    \item \textbf{H -- Architecture Change Management:} Änderungen
          am laufenden System kontrolliert einführen und den
          Zyklus bei Bedarf neu starten
\end{itemize}

Der Zyklus ist \textbf{iterativ}, d.\,h. nach Phase H kann der Prozess
wieder bei A (oder sogar Preliminary) beginnen, falls sich Anforderungen
grundlegend ändern.

\subsection{C4}

Das C4-Modell ist ein praxisorientiertes Framework zur Visualisierung
von Software-Architektur. Der Name steht für die vier Abstraktionsebenen:
\textbf{Context, Container, Component und Code}. Die Idee dahinter ist,
dass verschiedene Zielgruppen unterschiedliche Detailtiefe brauchen,
ein Manager braucht nur den groben Überblick, ein Entwickler hingegen
die genaue Komponentenstruktur.
\begin{figure}[H]
    \centering
    \includegraphics[width=\textwidth]{c4.png}
    \caption{Kubernetes Features}
\end{figure}

Die vier Ebenen bauen aufeinander auf:
\begin{itemize}
    \item \textbf{System Context:} Zeigt das System als Ganzes
          und seine Beziehungen zu Nutzern und externen Systemen
    \item \textbf{Container:} Aufschlüsselung in deploybare Einheiten
          wie Web-Apps, Datenbanken oder Microservices
    \item \textbf{Component:} Interne Struktur eines Containers,
          z.\,B. einzelne Module oder Services
    \item \textbf{Code:} Klassendiagramme oder ähnliches auf
          Code-Ebene (wird selten gezeichnet, da IDEs das übernehmen)
\end{itemize}

Ein großer Vorteil des C4-Modells ist die \textbf{Werkzeugfreiheit}.
Es funktioniert mit PlantUML, draw.io oder sogar auf einem Whiteboard,
solange die vier Ebenen klar getrennt bleiben.